\section{不读光与选定光学结果是两种实验}
\label{chap:feqo-conditioning}

式~\eqref{eq:feqo-partial-trace-components} 对所有未记录的光末态求和。如果实验还读出光场，可以只保留特定结果对应的电子。先看投影测量：光场投影 \(P_y=P_y^2=P_y^\dagger\) 代表结果 \(y\)，未归一化的条件电子态为
\begin{equation}
\widetilde\rho_{e|y}=\operatorname{Tr}_\gamma
\bigl[(I_e\otimes P_y)\rho_{e\gamma}(I_e\otimes P_y)\bigr],\qquad
p_y=\operatorname{Tr}\widetilde\rho_{e|y},\quad
\rho_{e|y}=\widetilde\rho_{e|y}/p_y.
\label{eq:feqo-conditional-electron-state}
\end{equation}
\(p_y=0\) 的结果不会发生，也就没有对应的条件态。把所有互斥结果的未归一化态相加，恢复不读光的 \(\rho_e\)；归一化后再相加则没有意义，必须带各自概率权重。

令前节二维示例在 \(gt=\pi/4\) 时的联合态为 \((|0,1\rangle-\ii|1,0\rangle)/\sqrt2\)。不读光得到 \(\rho_e=(|0\rangle\langle0|+|1\rangle\langle1|)/2\)，纯度 \(1/2\)。若光探测器投影到 \(|+\rangle_\gamma=(|0\rangle+|1\rangle)/\sqrt2\)，得到该结果的概率为 \(1/2\)，归一化的电子条件态为 \((|0\rangle-\ii|1\rangle)/\sqrt2\)，纯度为一。另一正交光学结果给相反的电子相位；忽略结果标签时，两种相位相加又消去。这就是“光场带走哪条路径信息”的可计算版本。

不必要求光学路径态恰好正交，才能量化这件事。考虑归一化纯态
\[
|\Psi\rangle=\sqrt p\,|0\rangle_e|u\rangle_\gamma
+e^{\ii\phi}\sqrt{1-p}\,|1\rangle_e|v\rangle_\gamma,
\qquad 0\leq p\leq1,
\]
其中 \(\langle u|u\rangle=\langle v|v\rangle=1\)，但 \(s:=\langle v|u\rangle\) 可以在零与单位模之间。对光求偏迹得到 \((\rho_e)_{01}=\sqrt{p(1-p)}e^{-\ii\phi}s\)，并由 \(2\times2\) 矩阵直接相乘求得
\begin{equation}
\operatorname{Tr}\rho_e^2=p^2+(1-p)^2+2p(1-p)|s|^2.
\label{eq:feqo-which-path-overlap}
\end{equation}
当 \(|u\rangle=|v\rangle\) 时 \(|s|=1\)，光没有记下电子走哪条能量路径，电子仍纯；当二者正交时 \(s=0\)，路径标签可由光完全区分，电子相干元消失。\(p\) 接近零或一时纯度也接近一，但那是因为电子几乎只走一条路径，不能误说环境“没有取得信息”。这个公式把路径信息、电子相干和纯度之间的关系落实到同一个可计算的内积。

光测量未必是正交投影。若结果 \(y\) 由作用于光场的算符 \(K_y\) 表示，则正算符 \(E_y=K_y^\dagger K_y\) 决定概率，\(\sum_yE_y=I\) 保证所有结果概率之和为一。把式~\eqref{eq:feqo-conditional-electron-state} 中的两个 \(P_y\) 分别换成 \(K_y\) 与 \(K_y^\dagger\) 就得到条件更新。结果效应 \(E_y\) 只决定概率；若要预测测量后继续演化的态，还必须给出具体 \(K_y\)。

电子谱变宽本身不证明纠缠：随机强度的经典光束也能混合出宽谱。这里的纠缠结论使用了已知的联合纯态及其约化纯度；实验验证则需要足以排除相应经典混合模型的电子—光相关或条件测量数据。

\begin{exercise}
对上面的联合态，把光投影改为 \(|-\rangle=(|0\rangle-|1\rangle)/\sqrt2\)，求电子条件态，并以两个结果的概率加权恢复无条件密度算符。
\end{exercise}
